% Generated draft: requires mathematical and source review.
\documentclass[11pt]{article}
\usepackage[margin=25mm]{geometry}
\usepackage{fontspec}
\setmainfont{Latin Modern Roman}
\usepackage{amsmath,amssymb,mathtools,unicode-math}
\setmathfont{Latin Modern Math}
\usepackage{enumitem,needspace,xurl,hyperref}
\setlength{\parindent}{0pt}
\setlength{\parskip}{5pt}
\setlength{\emergencystretch}{3em}
\newcommand{\N}{\mathbb{N}}
\newcommand{\Z}{\mathbb{Z}}
\newcommand{\Q}{\mathbb{Q}}
\newcommand{\R}{\mathbb{R}}
\newcommand{\C}{\mathbb{C}}
\newcommand{\F}{\mathbb{F}}
\newcommand{\A}{\mathbb{A}}
\newcommand{\PP}{\mathbb{P}}
\newcommand{\E}{\mathbb{E}}
\newcommand{\eps}{\varepsilon}
\newcommand{\dd}{\,\mathrm{d}}
\begin{document}
\section*{UnsolvedMath ID: 800899}
\textbf{Source problem code:} catalogue entry
\subsection{Definitions and mathematical statement}
Let \(X\) be a compact complex manifold of complex dimension \(n\), and let \(\omega\) be its Hermitian \((1,1)\)-form; no assumption that \(d\omega=0\) is imposed. Write \(\operatorname{PSH}(X,\omega)\) for the upper semicontinuous quasi-plurisubharmonic functions \(\varphi\) satisfying
\[
\omega+dd^c\varphi\geq 0
\]
in the sense of currents. Here quasi-plurisubharmonic means locally the sum of a plurisubharmonic function and a smooth function. The normalization of \(dd^c\) must be fixed consistently; the supplied record does not specify it. For bounded \(\omega\)-plurisubharmonic functions, interpret \((\omega+dd^c\varphi)^n\) as the complex Monge--Amp\`ere measure defined using Bedford--Taylor products and the smooth form \(\omega\).

Define
\[
\operatorname{cap}(X,\omega)
:=\inf\left\{
\int_X(\omega+dd^c\varphi)^n
\;:\;
\varphi\in\operatorname{PSH}(X,\omega),\quad
0\leq\varphi\leq1\ \text{on }X
\right\}.
\]
The pointwise bounds are part of the admissibility condition.

\textbf{Question (source Question 30).} Relate \(\operatorname{cap}(X,\omega)\) to the geometry of \((X,\omega)\). In particular, can this quantity be computed in terms of geometric data, and how does it behave under deformations of \((X,\omega)\)?

The supplied question is an open-ended research problem: it specifies neither the geometric data sought nor the class or regularity of deformations. In particular, it does not formulate a specific continuity, semicontinuity, or invariance assertion.
\subsection{Short English statement}
For a compact Hermitian manifold, take the infimum of the total complex Monge--Amp\`ere mass over all \(\omega\)-plurisubharmonic functions with values in \([0,1]\). Describe this quantity geometrically, determine whether geometric data can compute it, and study its behavior under deformations. The source leaves the geometric data, deformation framework, and normalization of \(dd^c\) unspecified.
\subsection{Sources}
\begin{itemize}
\item Local UnsolvedMath registry, permanent ID \(800899\), source problem code \texttt{catalogue entry}; title: \emph{Open Problems in Pluripotential Theory --- Question 30 (Dinew--Ko\l{}odziej)}. The supplied registry labels the problem \emph{open}; this label is supplied curation, not an independently verified assessment.
\item Dinew, Guedj, Zeriahi, \emph{Open Problems in Pluripotential Theory} (2015), Question 30. Supplied source link: \url{https://arxiv.org/abs/1511.00705}. The record reports access on 2026-10-16 and extraction by \texttt{arxiv\_latex}; these are supplied provenance details, not a claim of independent source access or verification.
\item The supplied background attributes the question or its context to Dinew--Ko\l{}odziej. It gives no separate reference clarifying that attribution.
\end{itemize}
\end{document}
